problem:  
Let \( M^n \) be a compact, simply connected Riemannian manifold with **positive sectional curvature** admitting an **effective, isometric torus action** of rank  
\[
r = \left\lfloor \frac{n+1}{2}\right\rfloor - 1.
\]  
Then \(M/T^r\) is a finite‑dimensional **Alexandrov space** with curvature \(\ge 1\), whose boundary \(\partial(M/T^r)\) represents the union of orbits with non‑trivial isotropy.  The maximal‑rank case \(r = \lfloor (n+1)/2 \rfloor\) is classified by the Grove–Searle theorem, and the orbit space in that case is known to be an Alexandrov *simplex*.  Little is known, however, when the symmetry rank drops by one.  

**Question:**  
In the near‑maximal symmetry case above, does positive curvature still enforce a strong **Alexandrov convexity condition** on the orbit space?  More precisely, is it true that for every boundary face \(F \subset \partial(M/T^r)\), the **angle between adjacent faces**, measured in the Alexandrov sense, satisfies  
\[
\angle(F_i,F_j) < \pi/2,
\]  
and that the boundary traces of singular orbits remain **totally geodesic subspaces** of positive curvature?  
Equivalently, does dropping the torus rank by one allow the orbit space to develop *non‑acute dihedral angles* or *curvature‑flat ridges*, or does positive sectional curvature rigidly prohibit such degenerations?  

Why it is a "good" problem:  
This version places the question squarely in the **intrinsic Alexandrov geometry of orbit spaces**, rather than in any embedding picture.  It asks whether the *acute‑angle convexity*—a geometric trace of maximal symmetry in positive curvature—survives under minimal relaxation of the symmetry‑rank constraint.  The problem links fine‑scale **metric curvature geometry** (Alexandrov angle conditions), **group action theory** (torus isotropy stratification), and **global Riemannian structure** (positive curvature).  

If true, the result would reveal an unexpected quantitative rigidity: near‑maximal symmetry still forces orbit spaces to behave with acute‑angle convexity, suggesting that isotropy faces “repel” each other under positive curvature.  If false, any counterexample would expose genuinely new Alexandrov geometries compatible with positive curvature and large symmetry rank.  The problem is simple to state, logically coherent, and sits at a fertile intersection of curvature and symmetry theory—thus meeting the criteria of a deep, “good” research problem.